\SourcePage{240}
\subsection{路径的接合}\label{sec:16h}
令 $\mu_0,\mu_1$ 为两条匹配路径，只在一个共同端点相交。假设路径的定向满足 $\mu_0(1)=\mu_1(-1)$；这只是为方便起见，定向在下面的构造中不起作用。考虑将 $\mu_k$ 拼起来，再移到共同端点右侧所得到的路径 $\mu=\mu_1\#\mu_0$，见\sr{fig:16.4}{图~16.4}。这是一种不对称的运算，依赖于如何排列这两条路径。显然，$\mu$ 仍为匹配路径。我们要解释的结果是：相应的匹配圈 $\Sigma_\mu$ 由 $\Sigma_{\mu_0}$ 沿 $\Sigma_{\mu_1}$ 作一次 Dehn 扭转得到。
\SourceFigure{figure-16-4}{0.25753818}{16.4}
为准确陈述这一结论，须把标架和同痕考虑在内。每个 $\mu_k$ 分为一对消失路径 $\gamma_k^\pm$，因而有四个相应的消失圈。方便起见，假设 $\mu|[-1,-1/2]$ 在显然的重新参数化下与 $\mu_0|[-1,0]$ 一致，类似地，$\mu|[1/2,1]$ 与 $\mu_1|[0,1]$ 一致。接合路径 $(\mu|[-1/2,0])^{-1}\cdot\gamma_0^+$ 与 $(\mu|[0,1/2])\cdot\gamma_1^-$ 保持端点不动地同伦；由此在 $E_{\mu(0)}$ 中得到带标架同痕
\begin{equation}\label{eq:16.9}
h_{\mu|[-1/2,0]}(V_{\gamma_0^+})\simeq
h_{\mu|[0,1/2]}^{-1}(V_{\gamma_1^-}).
\end{equation}
选择定义匹配圈 $\Sigma_{\mu_k,I_k}$ 所需的同痕 $I_k$。结合\eref{16.9}，便得到 $\mu$ 两半 $\gamma^\pm$ 的消失圈之间的同痕 $I$：
\[
V_{\gamma^-}=h_{\mu|[-1/2,0]}(V_{\gamma_0^-})
\overset I\simeq h_{\mu|[0,1/2]}^{-1}(V_{\gamma_1^+})=V_{\gamma^+},
\]
并用它定义另一个匹配圈 $\Sigma_{\mu,I}$。

\StatementNumber{13}\begin{lemma}\label{st:16.13}
在 $E$ 中存在带标架 Lagrangian 球面的同痕
\begin{equation}\label{eq:16.10}
\Sigma_{\mu,I}\simeq\tau_{\Sigma_{\mu_1,I_1}}(\Sigma_{\mu_0,I_0}).
\end{equation}
\end{lemma}
\begin{proof}
这是通过变形到更简单特殊情形来论证的典型例子。若结论对某个正合 Lefschetz 纤维化成立，则对其任意沿 $\partial^hE$ 行为良好的变形也成立；因为一旦为变形后的 Lefschetz 纤维化结构找到同痕，就可以
\SourcePage{241}
结合\stref{引理}{7.1}那样的一族共形辛嵌入，得到原结构的同痕。由\stref{引理}{15.3}及证明 Picard–Lefschetz 公式时的技巧，允许作下述假设：$\mu_0$ 两半对应的消失圈相等，$\mu_1$ 也如此；同痕 $I_0,I_1$ 恒定；且在 $E_{\mu_0(1)}$ 中唯一临界点附近，存在全纯坐标，使 $\pi$ 成为标准二次型\eref{15.9}，$\omega_E$ 成为 $\CC^{n+1}$ 上标准 Kähler 形式。仍不影响结论，还可变形所涉及的路径，使在 $\mu_0(1)$ 附近的局部坐标中，$\mu_0,\mu_1$ 分别为负、正实半轴，且 $\mu$ 在该点的小邻域之外与 $\mu_0\cup\mu_1$ 一致。于是，在 $E_{\mu_0(1)}$ 中临界点的小邻域之外，$\Sigma_{\mu,I}=\Sigma_{\mu_0,I_0}\cup\Sigma_{\mu_1,I_1}$；而在该点附近，由\stref{例题}{16.5}得到如下局部图景：
\begin{equation}\label{eq:16.11}
\Sigma_{\mu_0,I_0}=i\R^{n+1},\qquad
\Sigma_{\mu_1,I_1}=\R^{n+1},\qquad
\Sigma_{\mu,I}=\bigcup_t\sqrt{\mu(t)}S^n,
\end{equation}
其中 $\sqrt\mu$ 是 $\CC$ 中位于闭的右下象限内的光滑嵌入路径，满足 $t<-\epsilon$ 时 $\sqrt{\mu(t)}=it$，$t>\epsilon$ 时 $\sqrt{\mu(t)}=t$，$\epsilon$ 为某个小的正数。\znote{233}{局部模型\eref{16.11}是理解整个构造的关键，值得记牢：$\Sigma_{\mu_0}$ 在标准坐标里是虚方向的 $\R^{n+1}$，$\Sigma_{\mu_1}$ 是实方向的 $\R^{n+1}$，接合后的 $\Sigma_\mu=\bigcup_t\sqrt{\mu(t)}S^n$ 在 $t<-\epsilon$ 时落在 $i\R^{n+1}$ 中、$t>\epsilon$ 时落在 $\R^{n+1}$ 中，中间由半径随 $t$ 变化的球面插值。于是 $\Sigma_\mu$ 相当于把 $\Sigma_{\mu_0}$ 与 $\Sigma_{\mu_1}$ 各挖去一个开球，再沿两个半径 $\epsilon$ 的 $n$ 维球面用柄粘起来——这正是连通和，也是它与 $\Sigma_{\mu_1}\#\Sigma_{\mu_0}$ 同痕的直观原因。[来源：原书 §(16h)，印刷页 233。]}\eref{16.11}中第三个子流形，就是 Lagrangian 连通和或 Lagrangian 手术构造中出现的 \emph{Lagrangian 柄}。同样用 $\#$ 表示该连通和，并采用\bcite{125}小节 2e 的约定，有
\begin{equation}\label{eq:16.12}
\Sigma_{\mu,I}=\Sigma_{\mu_1,I_1}\#\Sigma_{\mu_0,I_0}.
\end{equation}
而\eref{16.10}与\eref{16.12}的右边在 Lagrangian 同痕意义下相同\bcite{124}，附录。同痕 $I_0,I_1,I$ 均恒定，所以容易验证这与标架相容。
\end{proof}

\StatementNumber{14}\begin{note}\label{st:16.14}
既然使用了记号 $\mu=\mu_1\#\mu_0$，忽略同伦后，\eref{16.10}可更直观地写为 $\Sigma_{\mu_1\#\mu_0}=\Sigma_{\mu_1}\#\Sigma_{\mu_0}$。还有另一种重要解释。令 $t_{\mu_1}$ 为沿 $\mu_1$ 的半扭转。由于 $\mu$ 与 $t_{\mu_1}(\mu_0)$ 同痕，我们证明的是
\[
\Sigma_{t_{\mu_1}(\mu_0)}\simeq\tau_{\Sigma_{\mu_1}}(\Sigma_{\mu_0}).
\]
在这种形式下，结论可推广到任意相交的匹配路径 $\mu_0,\mu_1$。但证明与上述特殊情形完全不同，此处也确无使用它的需要，故不作介绍；不过，\stref{引理}{18.3}可视为另一个特殊情形。
\end{note}

\subsection{将消失圈视为匹配圈}\label{sec:16i}
取正合 Lefschetz 双纤维化，记号如\eref{15.11}。第一个观察，是关于
\SourcePage{242}
临界值曲线 $\operatorname{Critv}(\varpi)\subset\mathcal S$ 及其到 $\mathcal W$ 的投影的拓扑的一个简单推论。对每个 $w\in\mathcal W\setminus(\operatorname{Critv}(\Psi)\cup\operatorname{Fakev}(\Psi,\varpi))$，考虑曲面 $\mathcal S_w$ 及其中的有限点构形 $\operatorname{Critv}(\varpi)\cap\mathcal S_w=\operatorname{Critv}(\varpi_w)$。由于交横截，该构形随 $w$ 光滑变化。沿任意路径 $\beta:[0,1]\to\mathcal W\setminus(\operatorname{Critv}(\Psi)\cup\operatorname{Fakev}(\Psi,\varpi))$ 追踪它的移动，就可定义\emph{辫平行移动映射}，即保定向微分同胚
\[
k_\beta:(\mathcal S_{\beta(0)},\operatorname{Critv}(\varpi_{\beta(0)}))
\longrightarrow(\mathcal S_{\beta(1)},\operatorname{Critv}(\varpi_{\beta(1)})),
\]
在同痕意义下唯一。可使它为辛，但由于维数低，这样做不会得到更多信息。现假设 $\gamma:[0,1]\to\mathcal W$ 为 $\Psi$ 的消失路径，且避开 $\operatorname{Fakev}(\Psi,\varpi)$。\stref{引理}{15.8}表明，当 $t\to1$ 时，$\operatorname{Critv}(\varpi)\cap\mathcal S_{\gamma(t)}$ 中的两个点，在覆盖 $\psi|\operatorname{Critv}(\varpi):\operatorname{Critv}(\varpi)\to\mathcal W$ 的一个分歧点处汇合。当 $s$ 接近 $1$ 时，取 $\mathcal S_{\gamma(s)}$ 中连接这两个点的短路径，它在不考虑定向的同痕意义下唯一。其在 $k_{\gamma|[0,s]}$ 下的原像，是嵌入路径 $\mu:[-1,1]\to\mathcal S_{\gamma(0)}$，满足 $\mu^{-1}(\operatorname{Critv}(\varpi_{\gamma(0)}))=\{-1,1\}$。称之为\emph{与 $\gamma$ 关联的匹配路径}，理由如下：

\StatementNumber{15}\begin{lemma}\label{st:16.15}
$\mu$ 是 $\varpi_{\gamma(0)}$ 的匹配路径，且适当选择同痕 $I$ 后，相应匹配圈 $\Sigma_{\mu,I}$ 在 $\mathcal E_{\gamma(0)}$ 中作为带标架 Lagrangian 球面，与消失圈 $V_\gamma$ 同痕。
\end{lemma}
\begin{proof}
只须证明：$s$ 接近 $1$ 时，在 $\mathcal E_{\gamma(s)}$ 中，连接 $\operatorname{Critv}(\varpi_{\gamma(s)})$ 中两个相近点的短路径是匹配路径，且所得匹配圈为 $\gamma|[s,1]$ 的消失圈。随后将一切从 $\mathcal E_{\gamma(s)}$ 移回 $\mathcal E_{\gamma(0)}$，就得到原结论。

考虑 $\mathcal E_{\gamma(1)}$ 中 $\Psi$ 的唯一临界点。在该点及其像附近引入具有\stref{引理}{15.9}所述性质的局部全纯坐标。如\eref{16.6}或\stref{引理}{16.13}的证明，可局部修改 $\mathcal E$ 上辛形式，使其在这些坐标中标准化。也可变形路径 $\gamma$，使在 $\Psi(x)\in\mathcal W$ 附近的局部坐标中，$t$ 接近 $1$ 时有 $\gamma(t)=1-t$。考虑 $s$ 接近 $1$ 时的短消失路径 $\gamma|[s,1]$。如\stref{例题}{16.5}，相应消失圈为
\[
V_{\gamma|[s,1]}=\sqrt{1-s}\,S^n\subset\CC^{n+1},
\]
另一方面，同一个球面也是 $\mathcal S_{\gamma(s)}$ 中连接 $\varpi_{\gamma(s)}$ 的两个临界值 $(y_1=1-s,y_2=\pm\sqrt{1-s})$ 的直线段所对应的匹配圈。验证这一点毫无困难，因为在此可用朴素程序\eref{16.8}定义匹配圈。最后，如\stref{引理}{16.13}的证明，应用\stref{引理}{7.1}回到原来的辛形式。当然，对该形式，匹配圈与消失圈不再严格相等，但论证给出两者之间的带标架同痕。\znote{234}{这条引理回答了第 16 节标题中“与”字的含义：消失圈与匹配圈在 Lefschetz 纤维化内部本是同一类对象，只是造法不同——消失圈沿 $\Psi$ 的一条消失路径生成，匹配圈则在曲面纤维 $\mathcal S_{\gamma(s)}$ 中取两个临界值之间的短路径、再经辫平行移动 $k_{\gamma|[0,s]}$ 拉回。注意上面的公式 $V_{\gamma|[s,1]}=\sqrt{1-s}\,S^n\subset\CC^{n+1}$ 中半径随 $s\to1$ 趋于零，正是“收缩到临界点”的显式写法；而同一条短直线段（$y_1=1-s$ 到 $y_2=\pm\sqrt{1-s}$）给出的匹配圈恰是同一个球面。[来源：原书 §(16i) 与引理 16.15，印刷页 234。]}
\end{proof}

\SourcePage{243}
\StatementNumber{16}\begin{note}\label{st:16.16}
$V_\gamma$ 的同痕类只依赖于 $\gamma$ 在全体消失路径空间中的同痕类；特别地，即使把 $\gamma$ 移过 $\operatorname{Fakev}(\Psi,\varpi)$ 的一个点，也不改变它。但辫平行移动通常会改变，因而关联的匹配路径 $\mu$ 也会改变。这并不矛盾，因为不同的匹配路径可以给出同一个匹配圈。事实上，\stref{例题}{16.12}中提到的这一现象，正是以这种方式从最简单的一般伪临界值产生的；在这些值处，曲线 $\operatorname{Critv}(\varpi)$ 具有尖点或结点。
\end{note}

\StatementNumber{17}\begin{note}\label{st:16.17}
如已提到，匹配圈由 Donaldson 引入，尚未发表。Lagrangian 连通和，尤其二维情形，研究于\bcite{109}。近来，它在镜像对称中也有了一定重要性，见\bcite{140}、\bcite{139}；后者包含公式\eref{16.12}的一个版本。熟悉辫单值变换技术，例如 Auroux 的分歧覆盖\bcite{11}与 Donaldson 束之间关系的读者，会觉得\secref{16i}的内容显而易见。
\end{note}
